\section{Method}

The pipeline has five stages. First, build or test commands are executed in a copied workspace. For heterogeneous external repositories, the build step can be explicitly skipped so that missing project-specific dependencies do not confound security-pipeline measurement. Second, scanners collect SAST, SCA, secret, Docker, and Kubernetes evidence. Third, findings are normalized into JSONL records with stable fingerprints and optional ground-truth links. Fourth, remediation rules patch only the copied workspace. Fifth, a second scanner pass measures residual findings before OPA evaluates the delivery policy.

Remediation is deliberately bounded. Deterministic rules update known vulnerable Python pins, remove documented fake secrets, escape reflected HTML output, parameterize a SQL query, and harden Docker or Kubernetes settings. For NodeGoat, the npm remediator runs \texttt{npm audit fix --package-lock-only --omit=dev} and records the npm return code and lockfile change. A non-zero npm return code is retained as evidence of residual advisories, not hidden as a tool failure.

OPA evaluates residual findings using a policy with zero tolerance for critical findings, a high-severity threshold of three, a CVSS ceiling of 9.0, and blocking on residual secrets. The policy counts residual findings as a list rather than a set so duplicate scanner findings are not silently deduplicated.
